\chapter{Εισαγωγή}\label{chap:intro}

Ο κύβος του Rubik αποτελεί κλασικό πρόβλημα αναζήτησης σε έναν μεγάλο διακριτό χώρο καταστάσεων. Το θέμα της εργασίας, όπως δόθηκε στο επίσημο κείμενο ανάθεσης, θέτει τρεις στόχους. Πρώτος στόχος είναι η υλοποίηση των αλγορίθμων Thistlethwaite, Kociemba και Korf σε γλώσσα Python ή/και Prolog και η εκτέλεσή τους σε συμβατικό υπολογιστή. Δεύτερος στόχος είναι η υλοποίηση αλγορίθμου που δέχεται μια κατάσταση και αναγνωρίζει πόσες κινήσεις απέχει από την τελική θέση. Τρίτος στόχος είναι η εύρεση κατάλληλης ευρετικής συνάρτησης για τη βέλτιστη επίλυση του κύβου με τον αλγόριθμο A* ή παραλλαγή του. Ως μία κίνηση μετρά η περιστροφή μιας έδρας κατά $\pm 90^\circ$ ή κατά $180^\circ$, σύμφωνα με τη μετρική μισών στροφών (\eng{half-turn metric}, HTM). Για τον τρίτο στόχο, η παρούσα εργασία επιλέγει ως παραλλαγή του A* τον αλγόριθμο IDA*, σε συνδυασμό με παραδεκτές ευρετικές που προκύπτουν από πίνακες προτύπων (\eng{pattern databases}, PDB), όπως αναλύεται στα Κεφάλαια~\ref{chap:algorithms} και~\ref{chap:implementation}.

Ως προς την επιλογή «Python ή/και Prolog», η παρούσα υλοποίηση επιλέγει την Python ως κύρια γλώσσα του συστήματος, για λόγους ελέγχου, δοκιμών και αναπαραγωγιμότητας· το επίσημο κείμενο ανάθεσης παραπέμπει επίσης στο σχετικό υλικό Prolog και Τεχνητής Νοημοσύνης του Κάλλιπου \cite{sgarbas2024prolog}. Οι κρίσιμοι για την απόδοση μηχανισμοί --- η παραγωγή μεγάλων πινάκων και οι βαθιές ακριβείς αναζητήσεις --- υλοποιούνται σε εγγενή κώδικα C/C++ που ενορχηστρώνεται από την Python. Η επιλογή αυτή υπαγορεύεται από τα όρια μνήμης και χρόνου εκτέλεσης ενός συμβατικού υπολογιστή με 16 GiB μνήμης, εντός των οποίων οι υπολογισμοί αυτοί δεν θα ήταν πρακτικά εφικτοί με αμιγώς διερμηνευόμενο κώδικα.

Ο γνωστός ισχυρισμός ότι κάθε επιλύσιμη κατάσταση λύνεται με το πολύ 20 κινήσεις στη μετρική μισών στροφών --- ο λεγόμενος αριθμός του Θεού (\eng{God's number}) --- έχει αποδειχθεί από τους Rokicki, Kociemba, Davidson και Dethridge \cite{rokicki2013diameter}, με τα συνοδευτικά δεδομένα του εγχειρήματος δημοσιευμένα στον σχετικό ιστότοπο \cite{rokicki2010godsnumber}. Το αποτέλεσμα αυτό χρησιμοποιείται εδώ ως εξωτερικό τεκμήριο και δεν αναπαράγεται ως πλήρης απόδειξη.

Κεντρική αρχή της εργασίας είναι η ακαδημαϊκή ακρίβεια των ισχυρισμών. Η λέξη «βέλτιστη» χρησιμοποιείται μόνο όταν μια πλήρης ή παραδεκτή αναζήτηση έχει ολοκληρωθεί για τη συγκεκριμένη κατάσταση· τα αποτελέσματα αυτά φέρουν την ετικέτα \code{exact} (αποδεδειγμένα βέλτιστη λύση). Αν μια διαδικασία επιστρέφει απλώς μια έγκυρη λύση, το αποτέλεσμα καταγράφεται ως \code{non_exact}, δηλαδή ως επαληθευμένη λύση χωρίς απόδειξη βελτιστότητας. Αν υπολογίζεται μόνο κάτω φράγμα, το αποτέλεσμα καταγράφεται ως \code{lower_bound} (αποδεδειγμένο κάτω φράγμα χωρίς λύση) και δεν παρουσιάζεται ως ακριβής απόσταση, ενώ η εξάντληση των ορίων αναζήτησης χωρίς λύση καταγράφεται ως \code{timeout}. Το λεξιλόγιο αυτό ορίζεται τυπικά στο Κεφάλαιο~\ref{chap:model} και χρησιμοποιείται με την ίδια σημασία σε ολόκληρη την εργασία.

\section{Κίνητρο}\label{sec:intro-motivation}
Ο κύβος του Rubik είναι ελκυστικό αντικείμενο για διπλωματική εργασία επειδή συνδυάζει τρεις διαφορετικές όψεις της πληροφορικής. Πρώτον, έχει καθαρή μαθηματική δομή: κάθε κίνηση μετασχηματίζει μια πεπερασμένη κατάσταση και όλες οι νόμιμες ακολουθίες κινήσεων μπορούν να θεωρηθούν στοιχεία ομάδας. Δεύτερον, έχει πρακτική αλγοριθμική δυσκολία: ο χώρος καταστάσεων του κύβου 3×3×3 αριθμεί 43\,252\,003\,274\,489\,856\,000 (περίπου $4{,}3\times 10^{19}$) προσβάσιμες καταστάσεις \cite{rokicki2013diameter}, ώστε η αφελής εξαντλητική αναζήτηση να γίνεται γρήγορα ανεφάρμοστη.\footnote{Σύμβαση αριθμητικής σημειογραφίας σε ολόκληρη την εργασία: οι ακέραιοι με πέντε ή περισσότερα ψηφία ομαδοποιούνται ανά τρία ψηφία με λεπτό τυπογραφικό διάστημα (π.χ.\ 42\,577\,920), ενώ οι δεκαδικές τιμές γράφονται με τελεία ως υποδιαστολή, ακριβώς όπως παράγονται από τα εργαλεία μετρήσεων (π.χ.\ 51.813928 δευτερόλεπτα).} Τρίτον, απαιτεί πειραματική πειθαρχία: ένας επιλύτης (\eng{solver}) μπορεί να βρίσκει λύσεις, αλλά αυτό δεν σημαίνει ότι αποδεικνύει ελάχιστο μήκος. Για τους ίδιους λόγους, το επίσημο κείμενο ανάθεσης αποδίδει στο θέμα υψηλό ερευνητικό χαρακτήρα.

Η διάκριση αυτή είναι ιδιαίτερα σημαντική σε μια ακαδημαϊκή εργασία με τίτλο που περιέχει τη λέξη «βέλτιστη». Στην καθημερινή χρήση η λέξη μπορεί να χρησιμοποιηθεί χαλαρά, για να δηλώσει μια πολύ καλή ή γρήγορη λύση. Στο πλαίσιο της παρούσας εργασίας, όμως, η βελτιστότητα είναι ισχυρισμός απόδειξης: για συγκεκριμένη κατάσταση σημαίνει ότι δεν υπάρχει συντομότερη λύση στην επιλεγμένη μετρική. Επομένως απαιτεί είτε πλήρη απαρίθμηση μέχρι το αντίστοιχο βάθος είτε παραδεκτή αναζήτηση που ολοκληρώνει όλα τα αναγκαία όρια. Όταν η αναζήτηση διακόπτεται από όρια χρόνου, βάθους ή κόμβων, το αποτέλεσμα καταγράφεται ως \code{timeout} ή ως κάτω φράγμα (\code{lower_bound}).

Η εργασία επιδιώκει να αναδείξει αυτή τη διάκριση με κώδικα και όχι μόνο με περιγραφή. Για τον κύβο 3×3×3 υλοποιούνται αναπαραστάσεις σε επίπεδο κυβιδίων (\eng{cubies}), έλεγχοι εγκυρότητας, διαδικασίες αναζήτησης και πίνακες προβολών. Παράγονται επίσης με εγγενή κώδικα ένας γωνιακός πίνακας προτύπων με $8!\cdot3^7 = 88\,179\,840$ καταστάσεις και οκτώ πίνακες προτύπων έξι ακμών. Για την ακριβή αναζήτηση πλήρους κύβου υπάρχει εγγενής επιλύτης IDA*, ο οποίος χρησιμοποιεί παραδεκτή ευρετική από τις τιμές αυτών των πινάκων και επιστρέφει \code{exact} μόνο όταν η αναζήτηση ολοκληρωθεί. Η αναγνώριση απόστασης αντιμετωπίζεται ως ξεχωριστό συμβόλαιο από την πρακτική επίλυση: επιστρέφει ακριβή απόσταση μόνο με απόδειξη, διαφορετικά κάτω φράγμα, αποτέλεσμα εξάντλησης ορίων ή ένδειξη άκυρης κατάστασης. Επειδή η πλήρης εξαντλητική απαρίθμηση ολόκληρου του χώρου 3×3×3 παραμένει εκτός των τοπικών υπολογιστικών πόρων, διατηρείται κανονικοποιημένη μελέτη του κύβου 2×2×2 (\eng{Pocket Cube}) ως υποστηρικτικό μικρό παράδειγμα. Με αυτόν τον τρόπο η παρούσα εργασία περιλαμβάνει πρακτικούς επιλύτες 3×3×3, πραγματική υποδομή πινάκων προτύπων για την IDA* και τεκμήρια ακριβούς αναζήτησης (\eng{exact-search evidence}) για συγκεκριμένες καταγεγραμμένες περιπτώσεις. Περιλαμβάνει επίσης μια εγγενή απόδειξη --- ανεξάρτητη από το εξωτερικό σύστημα H48/Nissy που περιγράφεται στο Κεφάλαιο~\ref{chap:implementation} --- ότι η απόσταση του \eng{superflip} είναι ακριβώς 20, καθώς και ένα πλήρες μικρότερο παράδειγμα όπου η εξαντλητική βέλτιστη αναζήτηση πράγματι ολοκληρώνεται.

\section{Ερευνητικά ερωτήματα}\label{sec:intro-questions}
Η εργασία οργανώνεται γύρω από τέσσερα ερωτήματα:

\begin{enumerate}
\item Πώς μπορεί να αναπαρασταθεί μια κατάσταση του κύβου με τρόπο που να επιτρέπει ελέγχους φυσικής εγκυρότητας, εφαρμογή κινήσεων και ανεξάρτητη επαλήθευση λύσεων;
\item Πώς διαφοροποιούνται στην πράξη η ρηχή εξαντλητική αναζήτηση, η IDA* με παραδεκτά κάτω φράγματα και οι πρακτικοί αλγόριθμοι τύπου Kociemba και Thistlethwaite;
\item Ποια αποτελέσματα μπορούν να τεκμηριωθούν από τοπικά παραγόμενα αρχεία μετρήσεων (\eng{benchmarks}) χωρίς να αποδοθούν ως καθολική βελτιστότητα;
\item Πώς συνδυάζονται πραγματικοί πίνακες προτύπων 3×3×3 με ένα μικρότερο πλήρες παράδειγμα 2×2×2 χωρίς να συγχέονται οι αντίστοιχοι ισχυρισμοί;
\end{enumerate}

Τα ερωτήματα αυτά συνδέουν τη θεωρία με την υλοποίηση. Η απάντηση στο πρώτο ερώτημα δίνεται μέσα από το μοντέλο κυβιδίων, τη συνάρτηση φυσικής εγκυρότητας και τον επαληθευτή λύσεων (\eng{verifier}). Η απάντηση στο δεύτερο δίνεται από τους διαφορετικούς επιλύτες και τις ετικέτες κατάστασης των αποτελεσμάτων. Η απάντηση στο τρίτο δίνεται από τα πρωτογενή (\eng{raw}) και τα επεξεργασμένα (\eng{processed}) αποτελέσματα του προφίλ εκτέλεσης \code{thesis}. Η απάντηση στο τέταρτο δίνεται από τον εγγενή γωνιακό πίνακα προτύπων, τους πίνακες προτύπων ακμών και τη μελέτη του κύβου 2×2×2. Στα αντίστοιχα αποτελέσματα διακρίνονται ρητά τα κάτω φράγματα 3×3×3, τα τεκμήρια ακριβούς αναζήτησης σε καταγεγραμμένες περιπτώσεις, η εγγενής απόδειξη απόστασης ακριβώς 20 για το \eng{superflip} και η πλήρης κατανομή αποστάσεων του 2×2×2.

\section{Συνεισφορές}\label{sec:intro-contributions}
Οι βασικές συνεισφορές της παρούσας εργασίας είναι οι εξής:

\begin{itemize}
\item υλοποίηση μοντέλου κυβιδίων με μεταθέσεις και προσανατολισμούς ακμών και γωνιών,
\item έλεγχος φυσικής εγκυρότητας και επιλυσιμότητας,
\item αναπαραγώγιμες διαταράξεις (\eng{scrambles}) και ανεξάρτητη επαλήθευση λύσεων,
\item εξαντλητική αναζήτηση BFS για ακριβείς ρηχές αποστάσεις,
\item αναζήτηση IDA* τύπου Korf με παραδεκτή ευρετική, παραγόμενους πίνακες προβολών, εγγενώς παραγόμενο γωνιακό πίνακα προτύπων 3×3×3 και οκτώ εγγενείς πίνακες προτύπων έξι ακμών,
\item εγγενής διαδρομή βέλτιστης αναζήτησης πλήρους κύβου, η οποία αποδεικνύει βελτιστότητα για επιλεγμένες περιπτώσεις 3×3×3 όταν η αναζήτηση ολοκληρώνεται,
\item εγγενής (χωρίς H48/Nissy) απόδειξη βελτιστότητας για το \eng{superflip}: συμμετρικά ανηγμένος πίνακας \eng{FlipUDSlice} φάσης 1 με 16 συμμετρίες, παραδεκτό κάτω φράγμα τριών αξόνων κατά Reid \cite{reid1997optimal} και εξαντλητικός IDA* ενιαίου ορίου (\eng{single-bound}) που αποκλείει κάθε λύση μήκους $\le 19$, ώστε μαζί με την επαληθευμένη λύση μήκους 20 να αποδεικνύεται απόσταση ακριβώς 20,
\item εγγενής οριοθετημένη (\eng{scoped}) διαδρομή Kociemba με πίνακες φάσης 1 και προβολές φάσης 2,
\item τετραφασική διαδρομή Thistlethwaite που εκτελεί τις μειώσεις υποομάδων ως άπληστη κάθοδο πάνω σε ακριβείς πίνακες αποστάσεων BFS, με εγγυημένο τερματισμό για κάθε έγκυρη κατάσταση,
\item πρακτικός προσαρμογέας της μεθόδου δύο φάσεων του Kociemba που επιστρέφει επαληθευμένες αλλά μη αποδεδειγμένα βέλτιστες λύσεις,
\item πλήρης κανονικοποιημένη μελέτη του κύβου 2×2×2 με αποθηκευμένη κατανομή αποστάσεων, ως υποστηρικτική μελέτη περίπτωσης με ακριβή αποτελέσματα,
\item αυτοματοποιημένα πειράματα, αποθηκευμένα αποτελέσματα και πίνακες που παράγονται από αυτά.
\end{itemize}

\section{Περιορισμός πεδίου}\label{sec:intro-scope}
Η εργασία δεν ισχυρίζεται ότι η υλοποίηση αποτελεί ταχύ μηχανισμό καθολικών απαντήσεων (\eng{oracle}) που λύνει βέλτιστα κάθε αυθαίρετη κατάσταση 3×3×3 σε πρακτικό τοπικό χρόνο. Ο εγγενής βέλτιστος επιλύτης δέχεται πλήρεις καταστάσεις σε επίπεδο κυβιδίων και δίνει βέλτιστη λύση όταν η παραδεκτή αναζήτηση ολοκληρώνεται· η πειραματική κάλυψη, ωστόσο, περιορίζεται στις συγκεκριμένες καταγεγραμμένες περιπτώσεις των αρχείων αποτελεσμάτων. Οι πλήρεις πίνακες αποκοπής (\eng{pruning tables}) και οι μεγάλοι πίνακες προτύπων που απαιτούνται για υψηλή καθολική απόδοση έχουν σημαντικό κόστος μνήμης και χρόνου, όπως φαίνεται από την κλασική εργασία του Korf \cite{korf1997pattern} και από την τεκμηρίωση της μεθόδου Kociemba \cite{kociembaImplementationDetails}.

Ο περιορισμός αυτός δεν είναι απλή τεχνική λεπτομέρεια· επηρεάζει τον τρόπο με τον οποίο διαβάζονται όλα τα αποτελέσματα. Όταν ένας επιλύτης επιστρέφει λύση με ετικέτα \code{non_exact}, η λύση έχει ελεγχθεί με εφαρμογή των κινήσεων στην αρχική κατάσταση, αλλά δεν έχει αποδειχθεί ότι είναι η συντομότερη δυνατή. Όταν εμφανίζεται \code{timeout}, η εργασία δεν συμπληρώνει το κενό με εικασία· η περίπτωση παραμένει αποτυπωμένη ως περιορισμός της συγκεκριμένης διαμόρφωσης. Αντίστοιχα, όταν η απόσταση αναγνωρίζεται μόνο ως κάτω φράγμα, το κείμενο δεν την παρουσιάζει ως ακριβή απόσταση.

Η επιλογή αυτή κάνει το αποτέλεσμα λιγότερο εντυπωσιακό ως προς τους ισχυρισμούς, αλλά περισσότερο χρήσιμο ως τεχνικό τεκμήριο. Ο αναγνώστης μπορεί να δει ποια μέρη είναι πλήρως ελεγμένα, ποια είναι πρακτικά επαληθευμένα και ποια παραμένουν εκτός τοπικού υπολογιστικού ορίου. Με αυτόν τον τρόπο η εργασία αποφεύγει το συνηθισμένο σφάλμα να συγχέεται ένας γρήγορος επιλύτης με έναν βέλτιστο επιλύτη.

\section{Δομή της εργασίας}\label{sec:intro-structure}
Η υπόλοιπη εργασία οργανώνεται ως εξής. Το Κεφάλαιο~\ref{chap:background} παρουσιάζει το θεωρητικό υπόβαθρο της αναζήτησης σε χώρους καταστάσεων, των παραδεκτών ευρετικών και των πινάκων προτύπων, και επισκοπεί τη σχετική βιβλιογραφία, από τον αλγόριθμο του Thistlethwaite έως τον αριθμό του Θεού. Το Κεφάλαιο~\ref{chap:model} ορίζει τη μοντελοποίηση του κύβου: την αναπαράσταση κυβιδίων, τους φυσικούς περιορισμούς επιλυσιμότητας, τη μετρική μισών στροφών, τις συντεταγμένες (\eng{coordinates}) που χρησιμοποιούνται για πίνακες και το λεξιλόγιο καταστάσεων των αποτελεσμάτων. Το Κεφάλαιο~\ref{chap:algorithms} αναλύει τους αλγορίθμους επίλυσης --- Thistlethwaite, Kociemba δύο φάσεων και Korf με IDA* --- και τον τρόπο με τον οποίο κάθε διαδρομή τεκμηριώνει τους ισχυρισμούς της. Το Κεφάλαιο~\ref{chap:design} περιγράφει την ανάλυση και τη σχεδίαση του συστήματος: την αρχιτεκτονική, τη ροή δεδομένων και τη διαδρομή παραγωγής των αποτελεσμάτων. Το Κεφάλαιο~\ref{chap:implementation} τεκμηριώνει την υλοποίηση, δηλαδή το πακέτο Python, τις εγγενείς μηχανές C/C++ και την παραγωγή των πινάκων προτύπων και αποκοπής. Το Κεφάλαιο~\ref{chap:usage} παρουσιάζει αναλυτικά παραδείγματα χρήσης του συστήματος, τα οποία δείχνουν βήμα προς βήμα τις βασικές λειτουργίες επίλυσης και αναγνώρισης απόστασης. Το Κεφάλαιο~\ref{chap:validation} περιγράφει την επαλήθευση των αποτελεσμάτων και την ακεραιότητα των ισχυρισμών, συμπεριλαμβανομένων των ανεξάρτητων διασταυρώσεων με εξωτερικά εργαλεία. Το Κεφάλαιο~\ref{chap:experiments} παρουσιάζει την πειραματική μεθοδολογία και τις μετρήσεις για όλες τις υλοποιημένες διαδρομές επίλυσης. Το Κεφάλαιο~\ref{chap:conclusions} συνοψίζει τη συζήτηση, τα συμπεράσματα, τους περιορισμούς και τις μελλοντικές κατευθύνσεις. Τέσσερα παραρτήματα συμπληρώνουν το κύριο σώμα: το Παράρτημα~\ref{app:repro} δίνει οδηγίες αναπαραγωγής των αποτελεσμάτων, το Παράρτημα~\ref{app:schema} τεκμηριώνει το σχήμα των αρχείων αποτελεσμάτων, το Παράρτημα~\ref{app:cli} αποτελεί πλήρη αναφορά της γραμμής εντολών και το Παράρτημα~\ref{app:ai} περιέχει την αναλυτική δήλωση χρήσης εργαλείων τεχνητής νοημοσύνης.
